\chapter{Conclusion and Future Work}
\label{ch:conclusion}

\section{Answers to the Research Questions}

\paragraph{RQ1.} With 256 labelled examples per task, routers raise the average score from 57.4 to $69.9\pm0.1$ and the gain over majority from 18.4\% to 59.1\%. A clear gap to the experts (78.2) remains.

\paragraph{RQ2.} The average rises from 60.8 at $k=16$ to 64.6 at $k=64$ and 69.9 at $k=256$, and it has not levelled off. Small budgets are unstable across seeds.

\paragraph{RQ3.} Fixed coefficients reach only 58.4, so better static coefficients alone do not explain the improvement, although the ablation does not separate input dependence from capacity.

\paragraph{RQ4.} Full-MLP and intermediate-only upscaling reach similar averages with different failures, on RTE and on MRPC respectively. Intermediate-only needs 47.2\% of the storage of five models, against 73.1\%. At batch size 32, the two are about 2.9 and 4.5 times slower than one BERT.

\section{Assessment of the Objectives}

Table~\ref{tab:goals} assesses the objectives of Section~\ref{sec:problem} against the evidence.

\begin{table}[htbp]
  \centering\tablestyle
  \caption{The objectives of the problem statement against the results.}
  \label{tab:goals}
  \begin{tabularx}{\linewidth}{>{\raggedright\arraybackslash}p{3.2cm} X >{\raggedright\arraybackslash}p{5.2cm}}
    \toprule
    Objective & Evidence & Assessment\\
    \midrule
    O1. Dynamic integration & Per-input MLP weights; fixed coefficients reach 58.4, against 69.9 & Met by design; the contribution of input dependence is supported but not isolated\\
    O2. Close to the task-specific models & 69.9 against 78.2, gain 59.1\%; RTE and MRPC remain weak; no task reaches its expert & Partly met\\
    O3. Fewer stored parameters & 73.1\% (full MLP) or 47.2\% (intermediate-only) of five models & Met, at a cost in latency\\
    \bottomrule
  \end{tabularx}
\end{table}

\section{Main Contributions and Practical Implications}

The contribution is practical, and no novel algorithm is claimed. This study adapts WEMoE to the post-norm BERT encoder with attention-mask-aware routing, applies a protocol that separates the data of each stage, and evaluates two upscaling scopes, three labelled budgets and an input-independent ablation against a majority-class floor. In practice, input-conditioned merging clearly beats the static merge on average but not the task-specific experts, and its case rests mainly on storage constraints.

\section{Future Work}

Merging four tasks without RTE would test whether RTE interferes with the other tasks. Larger budgets ($k=512$ and $1024$) would show whether the score levels off, and more FSFT steps would test the intermediate-only variant. Measuring task-vector norms per module type would check the reason for upscaling the MLP on BERT. A strict protocol would cap dev at $k$, and more seeds for the experts and the data split would give more reliable estimates. Per-example analysis would show what routing variation within a task follows.
